

\section{Introduzione}
\begin{frame}
 \frametitle{Introduzione al progetto di tesi}
 
 \begin{center}
\scalebox{.6}{\smartdiagram[circular diagram:clockwise]{Full core simulation, Analysis Thermal Hydraulics, Nuclear Database}}
\end{center}


\end{frame}

\section{Piattaforma numerica}
\begin{frame}
% 
% \frametitle{Introduction of multiphysics approach}
% Development of computational platform that integrates different computational tools, into the common wrapping framework given by the \textcolor{red}{\textbf{SALOME platform}}:
% \begin{itemize}
% \item Lattice code, \textcolor{blue}{\textbf{DRAGON}}
% \item Full core simulation code, \textcolor{bleudefrance}{\textbf{DONJON}}
% \item 3D-porous thermal-hydraulic CFD code, \textcolor{mygreen}{\textbf{FEMus}}
% \end{itemize}
% \begin{minipage}[c]{0.48\linewidth}
% \begin{center}
% \scalebox{.5}{\smartdiagram[bubble diagram]{
% Salome Platform,Lattice code,Core simulation code,CFD code
% }
% }
% \end{center}
% \end{minipage}
% \begin{minipage}[c]{0.48\linewidth}
% \begin{center}
% \scalebox{.6}{\smartdiagram[circular diagram:clockwise]{Full core simulation, Analysis Thermal Hydraulics, Nuclear Database}}
% \end{center}
% \end{minipage}\hfill
% 
% 
% \end{frame}
\section{Risultati}



% % % \section{introduction of approach}
% % % 
% % \begin{frame}
% % \frametitle{Introduction of multiphysics approach}
% % Development of computational platform that integrates different computational tools, into the common wrapping framework given by the \textcolor{red}{\textbf{SALOME platform}}:
% % \begin{itemize}
% % \item Lattice code, \textcolor{blue}{\textbf{DRAGON}}
% % \item Full core simulation code, \textcolor{bleudefrance}{\textbf{DONJON}}
% % \item 3D-porous thermal-hydraulic CFD code, \textcolor{mygreen}{\textbf{FEMus}}
% % \end{itemize}
% % \begin{minipage}[c]{0.48\linewidth}
% % \begin{center}
% % \scalebox{.5}{\smartdiagram[bubble diagram]{
% % Salome Platform,Lattice code,Core simulation code,CFD code
% % }
% % }
% % \end{center}
% % \end{minipage}
% % \begin{minipage}[c]{0.48\linewidth}
% % % \begin{center}
% % % \scalebox{.6}{\smartdiagram[circular diagram:clockwise]{Full core simulation, Analysis Thermal Hydraulics, Nuclear Database}}
% % % \end{center}
% % \end{minipage}\hfill


\end{frame}
% 
% 
% 
% 
% \section{target of approach}
% \begin{frame}[t]
% \frametitle{Targets of multiphysics approach}
% \begin{itemize}
% \item Investigate the coupling between neutronic code and thermal-hydraulic code
% \item Implementation of multiphysics approach into preliminary model of a Fast Small Modular Reactor (SMR) 
% \item Choosing a lead cooled reactor conceptual design, ALFRED, characterized by small core dimension along the axis
% \end{itemize}
% \begin{block}{Future perspectives}
% A coupled neutronics and thermal hydraulics analysis to better supporting the fast core dynamics (in both fields) at design level
% \end{block}
% \end{frame}
% 
% \section{description of neutronic code}
% \begin{frame}[t]
% \frametitle{Neutronic Code}
% The neutronic code is an \textcolor{blue}{open-source reactor physics code \texttt{"Dragon-Donjon Ver5.0.5"}} developed at Ecole Polytechnique de Montreal (Canada), designed around methods to solving transport equation of neutron.
% \begin{block}{Lattice code, DRAGON}
% Simulate the neutronic behavior of the fuel assembly in the reactor core and create the proper nuclear database, called MULTICOMPO
% \end{block}
% \begin{block}{Full core simulation, DONJON}
% Interpolate information contained in MULTICOMPO, simulate the comprehensive neutronic behavior of core and compute power and flux distribution
% \end{block}
% \end{frame}
% 
% \section{cross section}
% \begin{frame}[t]
% \frametitle{Nuclear Database - MULTICOMPO}
% \begin{block}{MULTICOMPO}
% \begin{itemize}
% 
% \item \textcolor{red}{\textbf{data structure}} which contain nuclear properties of assembly's materials.
% 
% \item possibility to recovery data concerned parameters as \textcolor{mygreen}{\textbf{nuclide's temperature and density or depletion calculation}}. 
% 
% \item necessity to create an optimal nuclear database for a correctly use of full core simulation code.
% 
% \end{itemize}
% \end{block}
% \begin{block}
% 
% \begin{enumerate}
% \item Isotopic cross-section library \texttt{JEFF 3.1.2} with 315 energy groups
% \item Lattice calculation
% \begin{itemize}
% \item Burn-up calculation
% \item Temperature and density modification
% \item Homogenization and collapse to 33 energy groups
% \end{itemize}
% \item Saving macroscopic cross sections in MULTICOMPO
% \end{enumerate}
% 
% \end{block}
% 
% \end{frame}
% 
% 
% \section{description of computational fluid dynamics code}
% \begin{frame}[t]
% \frametitle{Computational fluid dynamics code}
% % The 3D-CFD thermal hydraulics is based on the open-source code:
%  \begin{block}{FEMus}
% is a an open-source version, developed at the University of Bologna, of the LIBMESH library with parallel MPI-PETSC and multigrid solvers
% 
%  \end{block}
%  \begin{exampleblock}{The CFD solvers}
% have the function to \textcolor{blue}{solve energy and temperature equations} starting from a given \textcolor{blue}{neutronic power density distribution}
%  \end{exampleblock}
% 
% % 
% % \end{frame}
% % 
% % \begin{frame}
% % \frametitle{Temperature Field}
% \begin{minipage}{0.35\linewidth}
% \includegraphics[scale=0.15]{../img/tempeq.png}
% % Reactor
% \end{minipage}
% \begin{minipage}{0.28\linewidth}
% {\small
%  Assumption for solution of energy equation is \textcolor{blue}{porous medium}
%  \begin{equation}
%  \vec{v}\cdot\nabla T + \nabla K_{eff} \nabla T = Q_{source} \nonumber 
%  \end{equation}
% %  where $K_{eff}$ is thermal conductivity effective of the material
% }
% \end{minipage}
% \begin{minipage}{0.32\linewidth}
% \includegraphics[scale=0.2]{../img/canale.png}\\
% {\tiny Reactor channel (closed core)}
% \end{minipage}
% \end{frame}
% 
% 
% \begin{frame}
% 
% \frametitle{Computational platform}
% \framesubtitle{Open source codes}
% \begin{center}
% 
% \includegraphics[scale=0.40]{../img/platform.png}
% 
% \end{center}
% \end{frame}
% 
% \begin{frame}
% 
% \frametitle{Computational platform}
% \framesubtitle{Proprietary codes}
% \begin{center}
% 
% \includegraphics[scale=0.40]{../img/platform1.png}
% 
% \end{center}
% \end{frame}
% 
% \section{description of numerical platform}
% \begin{frame}
% \frametitle{Computational wrapping platform}
% \begin{alertblock}{The main purpose} 
% is to couple the solution of a multi-physics and multi-scale three-dimensional problem inside a simplified and more comprehensive framework.
% \end{alertblock}
% \begin{block}{The Salomé project}
% facilitates the coupling of scientific mesh-based codes thanks to its architecture and suite of tools that provide several data interfaces and exchange across the different codes
% \end{block}
% The integration of a code on the Salomè platform is obtained by generating an interface with functions available in the MEDMem library that \textcolor{red}{allows a data transfer from the platform to the code and then from the code to the interface}
% \end{frame}
% 
% \begin{frame}
% \frametitle{Main modules used in wrapping}
% \begin{block}{MED module}
% provide a library for \textcolor{blue}{\textbf{storing}} and \textcolor{blue}{\textbf{recovering}} computer data in a suitable format
% \end{block}
% \begin{block}{GEOM and MESH modules}
% have the main function to \textcolor{blue}{\textbf{draw}} and \textcolor{blue}{\textbf{create meshes}} in MED format, respectively.
% \end{block}
% \begin{block}{YACS module}
% is a tool to \textcolor{blue}{\textbf{supervise execution}} of complex interconnected scientific applications as object structure available during execution of DONJON code.
% \end{block}
% \end{frame}
% 
% \section{flowchart}
% \begin{frame}
% \frametitle{Summary flowchart}
% \begin{figure}
% \includegraphics[scale=0.30]{../img/flowchart_2.png}
% \end{figure}
% \end{frame}
% 
% \section{data modeling}
% 
% % \begin{frame}
% % \includegraphics[scale=0.2]{../img/crosection.pdf}
% 
% 
% % \end{frame}
% 
% 
% 
% 
% 
% 
% \begin{frame}
% \frametitle{The ALFRED (SMR) concept model}
% \framesubtitle{Data and Modeling}
% 
% \begin{minipage}{0.48\linewidth}
% \scalebox{.6}{
% \begin{tabular}{lcc}
% \hline
% Parameters & Unit & Value \\
% \hline
% Thermal Power & $MW$  & 300 \\
% Total height vessel & $mm$ & 3500 \\
% Inner vessel inner/outer radius & $mm$ & 1475/1525 \\ 
% Lattice pitch (hexagonal) & $mm$ & 13.86\\
% Wrapper thickness & $mm$ & 4.0\\
% Distance between to wrappers & $mm$ & 5.0\\
% Pins per FA & - & 126\\
% Total FA & - & 134\\
% Inner/Outer FA & - &  56/78 \\
% Inner/Outer enrichment & $\%$ & 20.5/26.2 \\
% Control/Safety Rods & - & 12/4 \\
% In-pile section & - & 1 \\
% Average Core Flow & $m/s$ & 1.28 \\
% Coolant Inlet/Outlet Temperature & $^oC$ & 400/520 \\
% \hline
% \end{tabular}
% }
% \begin{center}
% \includegraphics[scale=0.18]{../img/hexagon.pdf}
% \end{center}
% \end{minipage}\hfill
% \begin{minipage}{0.48\linewidth}
% \begin{center}
% Sketch of the reactor core
% \begin{figure}
% \includegraphics[scale=0.18]{../img/rectangle.pdf}
% \end{figure}
% \end{center}
% \end{minipage}
%   Reactor core fuel map 
% \end{frame}
% 
% 
% \begin{frame}
% \frametitle{Fuel pin and assembly}
% \framesubtitle{Data and Modeling}
% 
% 
% \scalebox{.7}{
% \begin{tabular}{lcc}
% \hline
% Parameters & Unit & Value \\
% \hline
% Pellet hollow diameter & $mm$ & 2.0\\
% Pellet radius & $mm$ & 4.5\\ 
% Gap thickness & $mm$ & 0.15\\
% Clad thickness & $mm$ & 0.6\\
% Pin diameter & $mm$ & 10.5\\
% Bottom plug length & $mm$ & 50\\
% Gas plenum height & $mm$ & 550\\
% Bottom plug length & $mm$ & 10\\
% Active height & $mm$ & 600\\
% Upper insulator height & $mm$ & 10\\
% Spring length & $mm$ & 120\\
% Upper plug length & $mm$ & 50\\
% \hline
% \end{tabular}
% % Dimensional Parameters of Fuel Pin
% }
% \begin{minipage}{0.3\linewidth}
% \begin{center}
% \includegraphics[scale=0.09]{../img/fullassembly.pdf}
% \end{center}
% $ \qquad $Fuel Assembly
% \end{minipage}
% \begin{minipage}[c]{0.2\linewidth}
% \begin{center}
% \includegraphics[scale=0.15]{../img/pin.pdf}
% \end{center}
% $ \qquad $Fuel pin
% \end{minipage}
% 
% 
% \end{frame}
% 
% 
% 
% 
% 
% \section{result}
% \begin{frame}[c]
% \frametitle{Finite elements for neutronic and CFD codes}
% \begin{minipage}[c]{0.48\linewidth}
% \centering Mesh CFD
% \begin{figure}
% \includegraphics[scale=0.3]{../img/1-ok.png}
% \end{figure}
% \end{minipage}
% \begin{minipage}[c]{0.48\linewidth}
% \centering Mesh Neutronic
% \begin{figure}
% \includegraphics[scale=0.3]{../img/2-ok.png}
% \end{figure}
% \end{minipage}
% \end{frame}
% \begin{frame}
% \frametitle{Preliminary result of multiphysics approach}
% \begin{minipage}[c]{0.32\linewidth}
% Flux distribution
% \begin{figure}
% \includegraphics[scale=0.14]{../img/a1.png}
% 
% \includegraphics[scale=0.14]{../img/b1.png}
% 
% \includegraphics[scale=0.14]{../img/c1.png}
% \end{figure}
% \end{minipage}
% \begin{minipage}[c]{0.32\linewidth}
% Power distribution
% \begin{figure}
% \includegraphics[scale=0.14]{../img/a2.png}
% 
% \includegraphics[scale=0.14]{../img/b2.png}
% 
% \includegraphics[scale=0.14]{../img/c2.png}
% \end{figure}
% \end{minipage}
% \begin{minipage}[c]{0.32\linewidth}
% Temperature field
% \begin{figure}
% \includegraphics[scale=0.14]{../img/a3.png}
% 
% \includegraphics[scale=0.14]{../img/b3.png}
% 
% \includegraphics[scale=0.14]{../img/c3.png}
% \end{figure}
% \end{minipage}
% \end{frame}
% 
% 
% 
% 
% \begin{frame}[c]
% \frametitle{Preliminary result of multiphysics approach}
% \begin{minipage}[c]{0.48\linewidth}
% \centering Power distribution along axial
% \begin{figure}
% \includegraphics[scale=0.25]{../img/pw.png}
% \end{figure}
% \end{minipage}
% \begin{minipage}[c]{0.48\linewidth}
% \centering Wappring power distribution with temperature field on reactor section
% \begin{figure}
% \includegraphics[scale=0.25]{../img/clipthre.png}
% \end{figure}
% \end{minipage}
% \end{frame}
% 
% \begin{frame}
% \frametitle{Conclusions}
% \begin{itemize}
% \item 
% \end{itemize}
% \end{frame}